\documentclass[11pt,a4paper]{article}
\usepackage[T1]{fontenc}
\usepackage[utf8]{inputenc}
\usepackage{lmodern}
\usepackage[margin=27mm]{geometry}
\usepackage{amsmath,amssymb,amsthm,bm}
\usepackage{booktabs,longtable,array}
\usepackage{url}
\usepackage{xurl}
\usepackage[hidelinks]{hyperref}
\usepackage{microtype}
\newcommand{\R}{\mathbb R}
\newcommand{\ran}{\operatorname{range}}
\newcommand{\rank}{\operatorname{rank}}
\newcommand{\col}{\operatorname{col}}
\newcommand{\norm}[1]{\left\lVert#1\right\rVert}
\newcommand{\code}[1]{\nolinkurl{#1}}
\newtheorem{proposition}{Proposition}[section]
\newtheorem{lemma}[proposition]{Lemma}
\newtheorem{example}[proposition]{Example}
\newtheorem{assumption}[proposition]{Assumption}
\theoremstyle{remark}
\newtheorem{remark}[proposition]{Remark}
\newtheorem{decision}[proposition]{Decision}

\title{Orthogonal-by-construction augmentation\newline
at the level of the state equation,\newline
for an augmentation carrying additional states}
\author{Master's thesis working derivation\\Eindhoven University of Technology}
\date{11 September 2026}

\begin{document}
\maketitle

\begin{abstract}
Gy\"or\"ok et al.\ (2026) parameterize an additive augmentation so that the learned component is
orthogonal to the baseline by construction, with no trade-off weight, and prove that the
resulting parameter Jacobian is block structured. Their model class is discrete-time
input--output with a baseline linear in its parameters and a static learning component. This note
carries the construction to the level of the state equation, for a baseline that is rational in
its parameters and self-scheduled, and for a learning component that carries its own additional
states. Four results. The condition on the augmented state map reduces, with no modelling choice,
to a condition on the learned write into the physical rows alone, because the baseline contributes
an identically zero block on the additional-state rows. Orthogonality holds at any rank, so the
rank deficiency of our regressor costs the uniqueness of the auxiliary coefficient and nothing
else. A coefficient frozen on an auxiliary point set makes the construction causal inside the
recursion and retains exactly the property the block-Jacobian argument needs. And, against the
construction, if that point set is the baseline-only rollout then the auxiliary coefficient is
provably independent of every weight on the latent path, so the projection charges nothing there;
separately, the one-step condition is not sufficient for non-overlap of the accumulated output.
Every step marked \emph{verified} was checked in exact symbolic arithmetic, on a minimal surrogate
and on the real gantry equations at explicit-Euler fidelity. Scripts and logs are named in
Section~\ref{sec:verif}.
\end{abstract}

\tableofcontents

\section{Scope, and what the level means}

Gy\"or\"ok's model class is discrete-time input--output: the regressor argument is a vector of
lagged measured inputs and outputs. Our plant model is a state-space LPV--LFR description, and
recasting it as an input--output model is the expensive path. The supervisors' direction is
therefore to impose the same orthogonality on the right-hand side of the state equation. The two
functions whose outputs are made orthogonal are the two state-update contributions that add up,
not the measured output. What the state-level condition does and does not say about the measured
output over a horizon is the subject of Section~\ref{sec:nonprop}, and it is a limitation of the
condition rather than an alternative reading of it.

Three items are out of scope and named so that they are not read into what follows: the value of
any penalty weight, since the construction has none; the existing fixed-reference trajectory
penalty, which is a verified method answering a different question and is not replaced here; and
the input--output level.

\paragraph{Notation.}
Physical state $x_p\in\R^{n_p}$ with $n_p=6$, additional states $x_a\in\R^{n_a}$ with $n_a=2$ by
default, $x=\col(x_p,x_a)\in\R^{n_x}$, input $u\in\R^{n_u}$ with $n_u=3$, output $y\in\R^{n_y}$
with $n_y=3$. Physical parameters $\theta\in\R^{n_\theta}$ with $n_\theta=14$, augmentation
parameters $\eta$, encoder parameters $\rho$. Selection and routing:
$E_p=\begin{bmatrix}I_{n_p}&0\end{bmatrix}^\top\in\R^{n_x\times n_p}$, and the routing expansion
splits as $R=\col(R_p,R_a)$ with $R_p\in\R^{n_p\times n_w}$ and $R_a\in\R^{n_a\times n_w}$. The
default configuration routes all $n_x$ rows, so $R_p=\begin{bmatrix}I_{n_p}&0\end{bmatrix}$ and
$R_a=\begin{bmatrix}0&I_{n_a}\end{bmatrix}$. The identity $E_p^\top R=R_p$ is used repeatedly.

\section{What is inherited}\label{sec:inherited}

For a baseline $\hat y_k=\phi(x_k)\theta_b$ linear in its parameters, and a static learning
component, stack over the estimation set and write $\Phi=\col_k\phi(x_k)$ and
$F_\eta=\col_k f^{\mathrm{ANN}}(x_k)$. The published construction is
\begin{align}
 \theta_{\mathrm{aux}}&=(\Phi^\top\Phi)^{-1}\Phi^\top F_\eta, \label{eq:g8}\\
 \widetilde F_\eta&=\bigl[I-\Phi(\Phi^\top\Phi)^{-1}\Phi^\top\bigr]F_\eta,\label{eq:g10}\\
 \Phi^\top\widetilde F_\eta&=0,\label{eq:g11}
\end{align}
which are Eqs.~(8), (10) and (11) of the paper, its Lemma~3 being \eqref{eq:g11}. At prediction
time the coefficient is frozen and the correction $\phi(x_k)\hat\theta_{\mathrm{aux}}$ is applied
at the current point, which is its Eq.~(13) and is what the released implementation does.

Two features of the source decide everything below, and both are easily missed.

\begin{remark}[The paper permits an auxiliary evaluation set]\label{rem:aux}
Immediately after Lemma~3 the paper states that $\theta_{\mathrm{aux}}$ may be built from any
auxiliary evaluation of the regressor, on a synthetically generated set or on a subset of the
estimation data, and that using the whole training set is a convenience rather than a restriction.
The freeze of Section~\ref{sec:construction} is therefore licensed by the source, not invented
here.
\end{remark}

\begin{remark}[What the block-Jacobian result actually needs]\label{rem:thm18}
The proof of the paper's Theorem~18 forms
$\psi_k=\begin{bmatrix}\phi(x_k)&J_f(x_k)\end{bmatrix}$ with
$J_f=\partial\widetilde f^{\mathrm{ANN}}/\partial\theta_a$ and argues
$\phi^\top J_f=0$, whence $\psi_k^\top\psi_k$ is block diagonal. That argument is
$J_f=(I-P_\Phi)\,\partial F_\eta/\partial\theta_a$ together with $\Phi^\top(I-P_\Phi)=0$, and it
is valid precisely because $P_\Phi$ does not depend on $\theta_a$. It does \emph{not} require
exactness against a regressor recomputed at the current iterate. A frozen projector is therefore
admissible for the block structure, and this is what makes Section~\ref{sec:construction}
possible.
\end{remark}

\subsection{Orthogonality holds at any rank}

Assumption~1 of the paper requires $\rank\Phi=n_\theta$. Ours fails that by construction. The
orthogonality statement survives regardless, and it is worth separating the two.

\begin{proposition}[Normal equations, any rank]\label{prop:anyrank}
Let $\Phi\in\R^{m\times n}$ and $G\in\R^{m}$ be arbitrary, and let $\lambda$ be any solution of
the normal equations $\Phi^\top\Phi\lambda=\Phi^\top G$. Then $\widetilde G=G-\Phi\lambda$
satisfies $\Phi^\top\widetilde G=0$.
\end{proposition}
\begin{proof}
$\Phi^\top\widetilde G=\Phi^\top G-\Phi^\top\Phi\lambda=0$ by the normal equations. A solution
exists for every $G$, since $\Phi^\top G\in\ran(\Phi^\top)=\ran(\Phi^\top\Phi)$.
\end{proof}

\begin{remark}
Rank deficiency makes $\lambda$ non-unique, and the pseudo-inverse of the released code selects
the minimum-norm solution. It does not repair the deficiency. What is lost with Assumption~1 is
the uniqueness argument of the paper's Theorem~7 and its Eq.~(21), not the orthogonality.
\emph{Verified}: Proposition~\ref{prop:anyrank} was checked on the surrogate at its true
deficient rank, and Eqs.~\eqref{eq:g8} to \eqref{eq:g11} were reproduced term for term on a
generic full-rank regressor.
\end{remark}

\section{The state equation, with the baseline separated}\label{sec:state}

The physical block reads $x_p$ and $u$ and writes the physical rows; the learned block reads the
\emph{full} state, additional states included, and writes the routed rows; the output block reads
$x_p$ and $u$ only. Hence
\begin{equation}\label{eq:model}
\boxed{
\begin{aligned}
 x_p(k+1)&=f_\theta\bigl(x_p(k),u(k)\bigr)\;+\;R_p\,g_\eta\bigl(x_p(k),x_a(k),u(k)\bigr),\\
 x_a(k+1)&=\phantom{f_\theta\bigl(x_p(k),u(k)\bigr)\;+\;}R_a\,g_\eta\bigl(x_p(k),x_a(k),u(k)\bigr),\\
 y(k)&=C\,x_p(k)+D\,u(k),\\
 x(0)&=e_\rho\bigl(y(-1),\dots,u(-1),\dots\bigr).
\end{aligned}}
\end{equation}
Four things to notice, each used later. The second row carries \emph{no} baseline term: the
baseline defines no equation for the additional states, so on that row there is nothing for a
learned contribution to overlap with. The learned block appears in both rows with the same
parameters $\eta$, so the two writes are not independent. The third row does not read $x_a$, so
the entire influence of the additional states on anything observable is routed through $R_pg_\eta$
at a later step. The fourth row contains parameters that appear nowhere else, so any condition
imposed on the first two rows is silent about the encoder.

Compactly, with $F_\theta(x,u)=E_pf_\theta(x_p,u)$ and $R=\col(R_p,R_a)$,
\begin{equation}\label{eq:compact}
 x(k+1)=F_\theta\bigl(x(k),u(k)\bigr)+R\,g_\eta\bigl(x(k),u(k)\bigr).
\end{equation}

\subsection{The baseline regressor and its offset}

Our baseline is rational in $\theta$ through $M(Y)^{-1}$ and then integrated, so $f_\theta$ is not
linear in $\theta$ and no exact $\phi$ exists. Expand at a fixed $\bar\theta$ with the evaluation
pair held fixed, which is the device the predecessor paper uses for nonlinear parameter dependence:
\begin{align}
 \Phi_p(x_p,u)&=\left.\frac{\partial f_\theta(x_p,u)}{\partial\theta}\right|_{\bar\theta}
   \in\R^{n_p\times n_\theta},
 &c(x_p,u)&=f_{\bar\theta}(x_p,u)-\Phi_p(x_p,u)\bar\theta\in\R^{n_p},\label{eq:taylor}\\
 \widetilde\Phi_p&=\begin{bmatrix}\Phi_p&c\end{bmatrix}\in\R^{n_p\times(n_\theta+1)}.
 &&\label{eq:extreg}
\end{align}

\begin{remark}[The offset column is not cosmetic]\label{rem:offset}
The supervisors' condition is that the learned output be orthogonal to \emph{the output of the
physics-based function}. For a baseline linear in its parameters, that set and the span of the
parameter sensitivities coincide, which is why the paper can use $\Phi$ directly. Ours is rational
in $\theta$, so they differ, and the faithful object is the affine set
$\ran\widetilde\Phi_p$, the nominal transition included. A contribution that cancels the
baseline's nominal response is charged under $\widetilde\Phi_p$ and is \emph{not} charged under
$\Phi_p$ alone. This is a deliberate departure from the earlier decision, recorded for the
trajectory penalty, to keep parameter-variation directions only; that method's declared target is
the parameter-displacement functional, this one's is non-overlap of the maps. \emph{Verified} on
the real gantry equations: the offset column is nonzero and lies outside the span of the fourteen
parameter columns, so appending it strictly enlarges the protected set. The price is that the
protected set now depends on the state origin, whereas parameter derivatives do not.
\end{remark}

\section{The orthogonality condition for this structure}\label{sec:condition}

\subsection{A per-step condition is unusable}

Stacking the full $n_x$ rows at a single step fails in both blocks at once. On the additional-state
rows the basis is identically zero, so nothing is constrained. On the physical rows,
$\widetilde\Phi_p$ has $n_\theta+1=15$ columns against $n_p=6$ rows, so it generically spans
$\R^{6}$, its orthogonal complement is $\{0\}$, and a per-step projection would force
$R_pg_\eta\equiv0$. The condition can only be a condition over a set of evaluation points. This is
why the source stacks.

\subsection{The condition, and why it sees only the physical rows}

Fix an evaluation set $\mathcal Z=\{(x_p(i),x_a(i),u(i))\}_{i=1}^{M}$ and stack the regressor of
the \emph{augmented} map \eqref{eq:compact} together with the learned write:
\begin{equation}\label{eq:stacks}
 \Psi=\col_i E_p\widetilde\Phi_p\bigl(x_p(i),u(i)\bigr)\in\R^{n_xM\times(n_\theta+1)},
 \qquad
 W_\eta=\col_i R\,g_\eta\bigl(x(i),u(i)\bigr)\in\R^{n_xM}.
\end{equation}

\begin{proposition}[The condition decouples onto the physical rows]\label{prop:decouple}
With $\Psi$ and $W_\eta$ as in \eqref{eq:stacks},
\begin{equation}\label{eq:decouple}
 \Psi^\top W_\eta=\sum_{i=1}^{M}\widetilde\Phi_p(i)^\top E_p^\top R\,g_\eta(i)
                 =\sum_{i=1}^{M}\widetilde\Phi_p(i)^\top R_p\,g_\eta(i),
\end{equation}
so the inner product does not involve $R_ag_\eta$ at all. Moreover any correction $\Psi\lambda$
writes into the physical rows only.
\end{proposition}
\begin{proof}
Immediate from $E_p^\top R=R_p$ and $\Psi\lambda=\col_iE_p\widetilde\Phi_p(i)\lambda$.
\end{proof}

\begin{remark}
This is the precise content of the informal statement that on the additional-state rows the
condition reads ``orthogonal to zero''. It is not a modelling choice to restrict the condition to
the physical write: the structure of \eqref{eq:model} does it. \emph{Verified} on the real gantry
equations: the parameter Jacobian of the augmented map is exactly zero on the additional-state
rows, as an identity in the state and the input.
\end{remark}

Accordingly, define the physical-row stacks
\begin{equation}\label{eq:phistack}
 \Phi_{\mathcal Z}=\col_i\widetilde\Phi_p\bigl(x_p(i),u(i)\bigr)\in\R^{n_pM\times(n_\theta+1)},
 \qquad
 G_\eta=\col_i R_p\,g_\eta\bigl(x(i),u(i)\bigr)\in\R^{n_pM},
\end{equation}
and state the condition.

\begin{decision}[State-level orthogonality condition]\label{dec:cond}
The augmentation is orthogonal to the baseline on $\mathcal Z$ when, for $\eta_{\mathrm{aux}}$
solving $\Phi_{\mathcal Z}^\top\Phi_{\mathcal Z}\eta_{\mathrm{aux}}=\Phi_{\mathcal Z}^\top G_\eta$,
\begin{equation}\label{eq:cond}
 \Phi_{\mathcal Z}^\top\widetilde G_\eta=0,
 \qquad \widetilde G_\eta=G_\eta-\Phi_{\mathcal Z}\eta_{\mathrm{aux}}.
\end{equation}
\end{decision}

By Proposition~\ref{prop:anyrank} this holds at any rank. By Proposition~\ref{prop:decouple} it
constrains the learned physical write only.

\section{The by-construction parameterization}\label{sec:construction}

\subsection{Why the coefficient cannot be fitted along the trajectory}

If $\mathcal Z$ is the training rollout then $x(i)$ depends on $\eta$, so
$\Phi_{\mathcal Z}$ does; worse, a coefficient fitted over a window depends on the learned writes
at later steps in that window, which depend on the coefficient. The correction at step $k$ would
be a function of the future. There is no exact fix inside one forward pass. The source never meets
this because its learning component is static.

\subsection{The construction}

Let $\mathcal Z_{\mathrm{ref}}$ be a point set fixed before the pass, and abbreviate
$\Phi_{\mathrm{ref}}=\Phi_{\mathcal Z_{\mathrm{ref}}}$ and
$K_{\mathrm{ref}}$ for any fixed left inverse solving the normal equations, both constant in
$\eta$ and $\theta$. Parameterize
\begin{equation}\label{eq:construction}
\boxed{
\begin{aligned}
 \eta_{\mathrm{aux}}(\eta)&=K_{\mathrm{ref}}\,
   \col_i R_p\,g_\eta\bigl(x_{\mathrm{ref}}(i),u_{\mathrm{ref}}(i)\bigr),\\
 x_p(k+1)&=f_\theta\bigl(x_p(k),u(k)\bigr)+R_p\,g_\eta\bigl(x(k),u(k)\bigr)
           -\widetilde\Phi_p\bigl(x_p(k),u(k)\bigr)\,\eta_{\mathrm{aux}}(\eta),\\
 x_a(k+1)&=R_a\,g_\eta\bigl(x(k),u(k)\bigr),\\
 y(k)&=C\,x_p(k)+D\,u(k).
\end{aligned}}
\end{equation}
The coefficient is a function of $\eta$ alone, computed once per pass, so the recursion is causal
and the map remains differentiable in $\eta$. The correction is evaluated at the \emph{current}
point and scaled by the \emph{frozen} coefficient, which is the form of the source's Eq.~(13).
There is no trade-off parameter.

\begin{proposition}[Exactness on the reference set]\label{prop:exact}
On $\mathcal Z_{\mathrm{ref}}$, the projected write of \eqref{eq:construction} satisfies
$\Phi_{\mathrm{ref}}^\top\widetilde G_\eta=0$ exactly, for every $\eta$ and at any rank.
\end{proposition}
\begin{proof}
Proposition~\ref{prop:anyrank} with $\Phi=\Phi_{\mathrm{ref}}$ and $G=G_\eta$ evaluated on
$\mathcal Z_{\mathrm{ref}}$.
\end{proof}

\begin{proposition}[The block structure survives the freeze]\label{prop:block}
Let $P_{\mathrm{ref}}=\Phi_{\mathrm{ref}}K_{\mathrm{ref}}$ be the associated projector. On
$\mathcal Z_{\mathrm{ref}}$,
\begin{equation}
 \frac{\partial}{\partial\eta}\Bigl[G_\eta-\Phi_{\mathrm{ref}}\eta_{\mathrm{aux}}(\eta)\Bigr]
 =(I-P_{\mathrm{ref}})\frac{\partial G_\eta}{\partial\eta},
 \qquad\text{hence}\qquad
 \Phi_{\mathrm{ref}}^\top\frac{\partial\widetilde G_\eta}{\partial\eta}=0 .
\end{equation}
\end{proposition}
\begin{proof}
$K_{\mathrm{ref}}$ is constant, so
$\partial(\Phi_{\mathrm{ref}}\eta_{\mathrm{aux}})/\partial\eta
=P_{\mathrm{ref}}\,\partial G_\eta/\partial\eta$; then
$\Phi_{\mathrm{ref}}^\top(I-P_{\mathrm{ref}})=0$ by Proposition~\ref{prop:anyrank} applied
columnwise.
\end{proof}

\begin{remark}
Proposition~\ref{prop:block} is the requirement identified in Remark~\ref{rem:thm18}, in our
notation. Two consequences for implementation. Differentiating \emph{through}
$\eta_{\mathrm{aux}}$ is not optional: treating it as a constant in the backward pass destroys the
property the construction exists to provide. And the guarantee is a statement about
$\mathcal Z_{\mathrm{ref}}$; along the training trajectory the residual
$\Phi_{\mathrm{traj}}^\top\widetilde G_\eta$ is not zero and its size is an empirical quantity of
the run. The honest claim is: orthogonal by construction with respect to the frozen reference set,
and approximately orthogonal elsewhere by an amount that must be measured.
\end{remark}

\section{What the additional states do}\label{sec:additional}

\begin{enumerate}
\item \textbf{They are the unconstrained part.} By Proposition~\ref{prop:decouple} the projector
acts as the identity on the $x_a$ block: no component of $R_ag_\eta$ is ever subtracted, at any
evaluation point, under any freeze.
\item \textbf{That is their intended role.} Those rows are where the learned component may carry
memory longer than one step, which is what the baseline provably cannot represent and therefore
what cannot be an overlap. Leaving them unconstrained is the correct design.
\item \textbf{Gauge invariance, as a correctness check.} Under $x_a\mapsto Vx_a$ with $V$
invertible, absorbed into the input and output weights of the learned block, $\Phi_{\mathrm{ref}}$
is unchanged and \eqref{eq:cond} is unchanged. The additional states are an arbitrary coordinate
system; a condition that constrained them would have to depend on that choice.
\item \textbf{Two things the construction does not do}, and they are the substance of
Sections~\ref{sec:zeroslice} and \ref{sec:nonprop}: the fitted coefficient can be blind to the
entire latent path, and the one-step condition does not control the accumulated output.
\end{enumerate}

\section{The zero slice: a defect of the obvious reference set}\label{sec:zeroslice}

The natural choice of $\mathcal Z_{\mathrm{ref}}$ is the closed-loop rollout of the baseline
alone, learned block gated off, which is the point set already specified elsewhere in this
project. In that rollout nothing writes the additional states, so $x_a\equiv0$ after the first
step. That choice is fatal to the part of the condition that matters most here.

\begin{proposition}[The frozen coefficient is blind to the latent path]\label{prop:zeroslice}
Suppose $x_a(i)=0$ for all $i$ in $\mathcal Z_{\mathrm{ref}}$. Let $\eta_{\mathrm{lat}}$ collect
every parameter of $g_\eta$ that acts only on the $x_a$ argument or only on the $R_a$ output
block. Then
\begin{equation}
 \frac{\partial\,\eta_{\mathrm{aux}}}{\partial\,\eta_{\mathrm{lat}}}=0
 \qquad\text{and}\qquad
 \frac{\partial\,G_\eta}{\partial\,\eta_{\mathrm{lat}}}=0 ,
\end{equation}
while the same parameters change the learned physical write at any trajectory point with
$x_a(k)\neq0$.
\end{proposition}
\begin{proof}
$G_\eta$ depends on $g_\eta$ only through $R_pg_\eta(x_p(i),0,u(i))$. Parameters acting on the
$x_a$ argument are multiplied by $x_a(i)=0$, and parameters acting on the $R_a$ block are
annihilated by $R_p$. Hence both derivatives vanish, and $\eta_{\mathrm{aux}}$, a fixed linear
map of $G_\eta$, inherits it. Off the reference set, $\partial(R_pg_\eta)/\partial\eta_{\mathrm{lat}}$
contains the factor $x_a(k)$.
\end{proof}

\emph{Verified} on the surrogate: all six latent-path parameters give an exactly zero derivative
of the fitted coefficient, while the physical write has derivative $x_a(k)$ with respect to the
latent-to-physical weight.

\begin{remark}[Consequence, with a verdict]
The projection subtracts nothing from the latent path, so the construction is not merely silent
about the $x_a$ rows, it is silent about the part of the \emph{physical} write that the additional
states produce. This invalidates any claim that the construction controls the augmentation's use
of its own states. It is a defect of the reference set, not of the condition.
\end{remark}

\subsection{Candidate repairs}

\begin{enumerate}
\item[(R1)] \textbf{Designed auxiliary set.} Draw $x_a$ over a designed region rather than from a
rollout, which Remark~\ref{rem:aux} permits explicitly. Causal and independent of $\eta$. Cost:
the points leave the data manifold, so the protected set describes a designed region, and the
region has to be justified.
\item[(R2)] \textbf{Lagged rollout.} Take $\mathcal Z_{\mathrm{ref}}$ from the augmented rollout
of the \emph{previous} epoch, refreshed between epochs. Within a pass the set is fixed, so
Propositions~\ref{prop:exact} and \ref{prop:block} hold exactly; the latent states are live, so
Proposition~\ref{prop:zeroslice} does not apply. Cost: the set depends on $\eta$ across epochs,
making the scheme a fixed-point iteration whose convergence is not established here. The
predecessor paper contemplates exactly this refresh for its projection matrix.
\item[(R3)] \textbf{Accept and declare.} Keep the baseline-only set and claim only non-overlap of
the instantaneous physical write at zero latent state.
\end{enumerate}

\begin{decision}[Recommended]
(R2), with (R1) as an independent check that the result does not depend on the refresh. (R2)
is the only candidate that keeps the evaluation points on the data manifold while making the
latent path visible, and it preserves both propositions within a pass, which is where the
block-structure argument lives.
\end{decision}

\section{One step is not enough for the output}\label{sec:nonprop}

Along a rollout of \eqref{eq:construction}, with $A(k)$ the state Jacobian of the closed augmented
map,
\begin{align}
 S(k+1)&=A(k)S(k)+E_p\Phi_p(k),
 &\frac{\partial y(k)}{\partial\theta}&=C\,E_p^\top S(k),\label{eq:Srec}\\
 T(k+1)&=A(k)T(k)+R\,\frac{\partial \widetilde g}{\partial\eta},
 &\frac{\partial y(k)}{\partial\eta}&=C\,E_p^\top T(k).\label{eq:Trec}
\end{align}
Condition \eqref{eq:cond} does not transfer through \eqref{eq:Srec} and \eqref{eq:Trec}, for three
independent reasons: the accumulation mixes time steps through $A(k)$, and orthogonality at equal
time says nothing about cross-time terms; $C$ retains $n_y=3$ of the $n_p=6$ physical rows, and
orthogonality in $\R^{6}$ is not preserved by a projection; and the additional-state write is
unconstrained yet enters $T$ through $R_a$ and returns to the physical rows through $A(k)$.

\begin{example}[One-step condition exact, output contribution not orthogonal]\label{ex:nonprop}
On the surrogate of Section~\ref{sec:verif}, take the sparsest live latent path: the position feeds
the additional state with weight $w_{21}$, the additional state feeds the physical velocity row
with weight $w_{13}$, all other weights zero. On the baseline-only reference set the learned write
is \emph{identically zero}, so \eqref{eq:cond} holds exactly and trivially. Rolling out six steps,
the learned contribution to the stacked output is nonzero from the fourth sample on, beginning
$\tfrac{1}{50}w_{13}w_{21}$ and
$\tfrac{112459}{1921000}w_{13}w_{21}$, and its inner product with the baseline output parameter
sensitivity is a nonzero rational function of $w_{13}w_{21}$. \emph{Verified} symbolically.
\end{example}

\begin{remark}[Claim wording]
The state-level condition is necessary and not sufficient for non-overlap of the accumulated
output. What it does buy is the removal of the one-step overlap, so the flat direction of the
source's Example~2, in which a learned linear layer trades off against the baseline parameters, is
eliminated at the level of the one-step augmented map. No wording may assert that the construction
prevents the augmentation from reproducing or negating baseline dynamics over a horizon.
\end{remark}

\section{Reduction to the published construction}\label{sec:reduction}

Three reductions are required, and the third is the level gap, named rather than hidden.

\begin{proposition}\label{prop:reduction}
Set $n_a=0$; route all physical rows, so $R_p=I_{n_p}$; let the baseline be linear in $\theta$, so
$c=0$ and $\widetilde\Phi_p=\Phi_p$ exactly; take $C=I$, $D=0$ with a one-step predictor and
measured states in place of the encoder; and take
$\mathcal Z_{\mathrm{ref}}$ to be the estimation set. Then \eqref{eq:model} becomes
$x(k+1)=\Phi_p(x(k),u(k))\theta+g_\eta(x(k),u(k))$, which is the source's Eq.~(5) under
$\phi\mapsto\Phi_p$ and $F_\eta\mapsto G_\eta$, and \eqref{eq:construction} together with
Decision~\ref{dec:cond} becomes \eqref{eq:g8}, \eqref{eq:g10} and \eqref{eq:g11} term for term,
with the inverse replaced by any normal-equations solution as in
Proposition~\ref{prop:anyrank}.
\end{proposition}

\begin{remark}
Without the third reduction the source's stacked object is the measured output and ours is the
state update, and those are not the same object. A derivation that claimed reduction without
naming it would be overstating the correspondence.
\end{remark}

\section{Assumptions of the source that our setting breaks}\label{sec:assumptions}

Each row carries a verdict, never a bare caveat.

\begin{longtable}{@{}p{0.30\linewidth}p{0.62\linewidth}@{}}
\toprule
\textbf{Assumption} & \textbf{Our setting, and the verdict}\\
\midrule
\endhead
Baseline linear in the parameters, Eq.~(2)
& Rational in $\theta$ through $M(Y)^{-1}$, then integrated. \emph{Weakens to} orthogonality
against the affine tangent set at $\bar\theta$. Theorem~7 does not transfer. The offset column of
Remark~\ref{rem:offset} restores the nominal response, at the price of dependence on the state
origin.\\
Regressor argument is a measured lagged input--output vector
& Ours are unmeasured states. This is the extension the source's conclusion names as open.
\emph{Weakens to} a condition relative to a chosen reference set, and introduces a dependence of
the regressor on $\bar\theta$ that the source does not have.\\
No encoder
& $\rho$ appears in no row the condition touches. \emph{Invalidates} coverage of a contribution
injected purely through $x(0)$.\\
Static learning component
& Ours carries $n_a$ states with no baseline equation. On that block the guarantee is
\emph{absent}, and with Section~\ref{sec:nonprop} this \emph{invalidates} horizon-level
non-overlap claims. With a baseline-only reference set, Section~\ref{sec:zeroslice} further
\emph{invalidates} control of the latent path until (R2) or (R1) is adopted.\\
No recursion
& \emph{Weakens} the property from ``with respect to the current trajectory'' to ``with respect to
the frozen reference set''. The block-structure argument survives, by
Proposition~\ref{prop:block}.\\
Assumption~1, $\rank\Phi=n_\theta$
& Fails structurally. \emph{Verified}: the parameter block has rank $10$ of $14$, and $11$ of $15$
with the offset column. \emph{Weakens to} uniqueness on $\ran\Phi^\top$ only; orthogonality
unaffected, by Proposition~\ref{prop:anyrank}.\\
Condition~4, $\Phi^\top\Delta=0$ on the data
& $\Delta$ is unknown outside simulation, so Theorem~7 cannot be invoked. \emph{Weakens to} the
error expression of Eq.~(21) plus the comparative statement under the source's Assumption~8.\\
$\Sigma_e$ diagonal, for Theorem~18
& Noiseless simulation has no noise term; closed-loop channels are not known to be uncorrelated.
\emph{Weakens to} block diagonality of $\Psi^\top\Psi$ without the covariance conclusion.\\
Conditions~11 and~14, stable predictor and persistency of excitation
& The augmented rollout is marginally stable on the two integrator axes. \emph{Invalidates} the
consistency results. The source's Remark~13 points at stable-by-design parameterizations.\\
A frozen basis transfers
& $\widetilde\Phi_p$ depends on the scheduling variable, and held-out scheduling positions are a
validation axis. \emph{Weakens to} exactness over the range the reference set visits; mitigation
is to span the operating range by design.\\
One-step prediction cost
& Ours is free-run with multiple shooting. \emph{Weakens} nothing in the condition, which never
referred to the cost, but removes comparability with the source's numerical results.\\
\bottomrule
\end{longtable}

\section{Verification record}\label{sec:verif}

All arithmetic below is exact rational or symbolic. Scripts and logs live in
\code{scripts/gantry/orthogonal-by-construction/code/}.

\paragraph{Surrogate, \code{surrogate\_obc\_symbolic.m}.}
Two physical states, one additional state, three parameters of which two enter only as a sum and
one enters rationally, self-scheduled on a state, routing to the velocity row and the latent row,
an output keeping the position and discarding the velocity, and a learning component affine in its
weights. Confirms: the exact column dependency and the resulting deficient rank;
Proposition~\ref{prop:anyrank} at that rank; Proposition~\ref{prop:reduction} against
\eqref{eq:g8} to \eqref{eq:g11}; Proposition~\ref{prop:zeroslice} for all six latent-path
parameters; and Example~\ref{ex:nonprop}.

\paragraph{Real equations at Euler fidelity, \code{gantry\_euler\_regressor\_structure.m}.}
The true $M(Y)=M_0+M_1Y+M_2Y^2$ with $M(Y)^{-1}$ kept rational, all fourteen parameters, one
explicit Euler step, self-scheduled on the third state. Explicit Euler replaces the integrator of
the implementation deliberately: it changes none of the structural properties checked, and keeps
the expressions tractable. Confirms Proposition~\ref{prop:decouple} on the real model, the three
exact column dependencies as identities in the state and the input, and the offset column being
nonzero and outside the parameter span.

\begin{remark}[An identifiability finding, to be reconciled]
The null space of the parameter block has four directions, not three. Besides the two expected
pairs it contains
\begin{equation}
 -\tfrac{1}{2}\bigl(\delta m_1+\delta m_2\bigr)+\delta m_b+\tfrac{L_b^2}{4}\,\delta J=0,
\end{equation}
for $J$ either inertia, whose difference recovers the expected inertia pair. So the inertia sum is
identifiable only in combination with the mass sum through
$J_b+J_h+\tfrac{L_b^2}{4}(m_1+m_2)$, which is exactly the entry of $M_0$ they share. The
project's documented set of ten identifiable quantities lists the inertia sum on its own. The two
statements concern different objects, a one-step map with the full state visible against a
trajectory read through the output matrix, so this is not yet a contradiction. It should be
reconciled before either is quoted in the thesis.
\end{remark}

\paragraph{Not verified here.} The size of the trajectory residual of
Proposition~\ref{prop:block}, which decides whether the frozen construction behaves like a
construction or like a weak penalty in practice; the achieved rank on real data rather than at
designed points; and everything downstream of a training run, since no implementation of
\eqref{eq:construction} exists.

\end{document}
